\chapter{2006年湖南卷（理）}
\section{单选题}
\begin{problem}
函数 \( y = \sqrt{\log_{2} x - 2} \) 的定义域是（\quad）

\choices
    {\((3, +\infty)\)}
    {\([3, +\infty)\)}
    {\((4, +\infty)\)}
    {\([4, +\infty)\)}
\end{problem}

\begin{answer}
D
\end{answer}

\begin{problem}
若数列 \(\{a_n\}\) 满足：\(a_{1} = \frac{1}{3}\)，且对任意正整数 \(m,n\) 都有 \(a_{m+n} = a_m \cdot a_n\)，则 \(\displaystyle\lim_{n \to +\infty} (a_1 + a_2 + \cdots + a_n) =\)（\quad）

\choices
    {\(\frac{1}{2}\)}
    {\(\frac{2}{3}\)}
    {\(\frac{3}{2}\)}
    {\(2\)}
\end{problem}

\begin{answer}
A
\end{answer}

\begin{problem}
过平行六面体 \(ABCD - A_{1}B_{1}C_{1}D_{1}\) 任意两条棱的中点作直线，其中与平面 \(DBB_{1}D_{1}\) 平行的直线共有（\quad）

\choices
    {\(4\) 条}
    {\(6\) 条}
    {\(8\) 条}
    {\(12\) 条}
\end{problem}

\begin{answer}
D
\end{answer}

\begin{problem}
“\(a=1\)”是“函数 \( f(x) = |x - a| \) 在区间 \( [1, +\infty) \) 上为增函数”的（\quad）

\choices
    {充分不必要条件}
    {必要不充分条件}
    {充要条件}
    {既不充分也不必要条件}
\end{problem}

\begin{answer}
A
\end{answer}

\begin{problem}
已知 \( \left|\bs{a}\right| = 2\left|\bs{b}\right| \neq 0 \)，且关于 \(x\) 的方程 \(x^{2} + \left|\bs{a}\right|x + \bs{a}\cdot\bs{b} = 0\) 有实根，则 \( \bs{a} \) 与 \( \bs{b} \) 的夹角的取值范围是（\quad）

\choices
    {\(\left[0, \frac{\pi}{6}\right]\)}
    {\(\left[\frac{\pi}{3}, \pi\right]\)}
    {\(\left[\frac{\pi}{3}, \frac{2\pi}{3}\right]\)}
    {\(\left[\frac{\pi}{6}, \pi\right]\)}
\end{problem}

\begin{answer}
B
\end{answer}

\begin{problem}
某外商计划在 \(4\) 个候选城市投资 \(3\) 个不同的项目，且在同一个城市投资的项目不超过 \(2\) 个，则该外商不同的投资方案有（\quad）

\choices
    {\(16\) 种}
    {\(36\) 种}
    {\(42\) 种}
    {\(60\) 种}
\end{problem}

\begin{answer}
D
\end{answer}

\begin{problem}
过双曲线 \(M: x^{2} - \frac{y^{2}}{b^{2}} = 1\) 的左顶点 \(A\) 作斜率为 \(1\) 的直线 \(l\)，若 \(l\) 与双曲线 \(M\) 的两条渐近线分别相交于点 \(B\)，\(C\)，且 \(|AB| = |BC|\)，则双曲线 \(M\) 的离心率是（\quad）

\choices
    {\(\sqrt{10}\)}
    {\(\sqrt{5}\)}
    {\(\frac{\sqrt{10}}{3}\)}
    {\(\frac{\sqrt{5}}{2}\)}
\end{problem}

\begin{answer}
A
\end{answer}

\begin{problem}
设函数 \( f(x) = \frac{x - a}{x - 1} \)，集合 \( M = \{x \mid f(x) < 0\} \)，\( P = \{x \mid f'(x) > 0\} \)，若 \( M \not\subset P \)，则实数 \( a \) 的取值范围是（\quad）

\choices
    {\((- \infty, 1)\)}
    {\((0,1)\)}
    {\((1, +\infty)\)}
    {\([1, +\infty)\)}
\end{problem}

\begin{answer}
A
\end{answer}

\begin{problem}
棱长为 \(2\) 的正四面体的四个顶点都在同一个球面上，若过该球球心的一个截面如图，则图中三角形（正四面体的截面）的面积是（\quad）

\FigureLayoutDeclare{img/2006/hunan_science/q9_fig1.png}{4.0cm}{9bp 13bp 7bp 8bp}
\bitmapfigure[max width=0.34\linewidth]{img/2006/hunan_science/q9_fig1.png}

\choices
    {\(\frac{\sqrt{2}}{2}\)}
    {\(\frac{\sqrt{3}}{2}\)}
    {\(\sqrt{2}\)}
    {\(\sqrt{3}\)}
\end{problem}

\begin{answer}
C
\end{answer}

\begin{problem}
若圆 \(x^{2} + y^{2} - 4x - 4y - 10 = 0\) 上至少有三个不同的点到直线 \(l: ax + by = 0\) 的距离为 \(2\sqrt{2}\)，则直线 \(l\) 的倾斜角的取值范围是（\quad）

\choices
    {\(\left[\frac{\pi}{12}, \frac{\pi}{4}\right]\)}
    {\(\left[\frac{\pi}{12}, \frac{5\pi}{12}\right]\)}
    {\(\left[\frac{\pi}{6}, \frac{\pi}{3}\right]\)}
    {\(\left[0, \frac{\pi}{2}\right]\)}
\end{problem}

\begin{answer}
B
\end{answer}

\section{填空题}
\begin{problem}
若 \((ax - 1)^5\) 的展开式中 \(x^3\) 的系数是 \(-80\)，则实数 \(a\) 的值是\fillinblank{}.
\end{problem}

\begin{answer}
\(-2\)
\end{answer}

\begin{problem}
已知 \(\begin{cases}
x \geqslant 1,\\
x - y + 1 \leqslant 0,\\
2x - y - 2 \leqslant 0
\end{cases}\)，则 \(x^2 + y^2\) 的最小值是\fillinblank{}.
\end{problem}

\begin{answer}
\(5\)
\end{answer}

\begin{problem}
曲线 \(y = \frac{1}{x}\) 和 \(y = x^2\) 在它们交点处的两条切线与 \(x\) 轴所围成的三角形的面积是\fillinblank{}.
\end{problem}

\begin{answer}
\(\frac{3}{4}\)
\end{answer}

\begin{problem}
若 \( f(x) = a \sin \left( x + \frac{\pi}{4} \right) + b \sin \left( x - \frac{\pi}{4} \right) \)（\(ab \neq 0\)）是偶函数，则有序实数对 \((a, b)\) 可以是\fillinblank{}. （写出你认为正确的一组数字即可）
\end{problem}

\begin{answer}
\((-1,1)\)
\end{answer}

\begin{problem}
如图，\(OM \parallel AB\)，点 \(P\) 在由射线 \(OM\)，线段 \(OB\) 及 \(AB\) 的延长线围成的区域内（不含边界）运动，且 \(\overrightarrow{OP} = x\overrightarrow{OA} + y\overrightarrow{OB}\)，则 \(x\) 的取值范围是\fillinblank{}；当 \(x = -\frac{1}{2}\) 时，\(y\) 的取值范围是\fillinblank{}.

\bitmapfigure[max width=0.55\linewidth]{img/2006/hunan_science/q15_fig1.png}
\end{problem}

\begin{answer}
\((-\infty,0)\)，\(\left(\frac{1}{2},\frac{3}{2}\right)\)
\end{answer}

\section{解答题}
\begin{problem}
如图，\(D\) 是直角 \(\triangle ABC\) 斜边 \(BC\) 上一点，\(AB = AD\)，记 \(\angle CAD = \alpha\)，\(\angle ABC = \beta\).

\begin{enumerate}
\item 证明：\(\sin \alpha + \cos 2\beta = 0\)；
\item 若 \(AC = \sqrt{3}DC\)，求 \(\beta\) 的值.
\end{enumerate}

\bitmapfigure[max width=0.34\linewidth]{img/2006/hunan_science/q16_fig1.png}
\end{problem}

\begin{solution}
\begin{enumerate}
\item 由于 \(D\) 在 \(BC\) 上，所以 \(\angle ABD=\angle ABC=\beta\). 又 \(AB=AD\)，故 \(\angle ADB=\angle ABD=\beta\)，从而 \(\angle BAD=\pi-2\beta\). 在直角三角形 \(ABC\) 中，\(\angle BAC=\dfrac{\pi}{2}\)，所以 \(\alpha+\pi-2\beta=\dfrac{\pi}{2}\)，即 \(\alpha=2\beta-\dfrac{\pi}{2}\). 因此 \(\sin\alpha=\sin\left(2\beta-\dfrac{\pi}{2}\right)=-\cos2\beta\)，所以 \(\sin\alpha+\cos2\beta=0\).

\item 在三角形 \(ADC\) 中，\(\angle ACD=\angle ACB=\dfrac{\pi}{2}-\beta\)，且由上一问知 \(\alpha=2\beta-\dfrac{\pi}{2}\)，故 \(\angle ADC=\pi-\beta\). 由正弦定理，
\[
\frac{AC}{DC}=\frac{\sin\angle ADC}{\sin\angle CAD}=\frac{\sin\beta}{\sin\alpha}=\sqrt{3}.
\]
由于 \(\alpha>0\)，有 \(\dfrac{\pi}{4}<\beta<\dfrac{\pi}{2}\). 令 \(s=\sin\beta\)，则
\[
\sin\alpha=\sin\left(2\beta-\frac{\pi}{2}\right)=-\cos2\beta=1-2s^{2}.
\]
所以 \(\sqrt{3}(1-2s^{2})=s\)，即 \(2\sqrt{3}s^{2}+s-\sqrt{3}=0\). 解得 \(s=\dfrac{1}{\sqrt{3}}\) 或 \(s=-\dfrac{\sqrt{3}}{2}\)，结合 \(0<s<1\) 可知 \(s=\dfrac{\sqrt{3}}{3}\). 因而 \(\beta=\arcsin\dfrac{\sqrt{3}}{3}\).
\end{enumerate}
\end{solution}

\begin{problem}
某安全生产监督部门对 \(5\) 家小型煤矿进行安全检查（简称安检）. 若安检不合格，则必须整改. 若整改后经复查仍不合格，则强制关闭. 设每家煤矿安检是否合格是相互独立的，且每家煤矿整改前安检合格的概率是 \(0.5\)，整改后安检合格的概率是 \(0.8\)，计算（结果精确到 \(0.01\)）：

\begin{enumerate}
\item 恰好有两家煤矿必须整改的概率；
\item 平均有多少家煤矿必须整改；
\item 至少关闭一家煤矿的概率.
\end{enumerate}
\end{problem}

\begin{solution}
\begin{enumerate}
\item 设必须整改的煤矿数为 \(X\). 每家煤矿必须整改的概率为 \(1-0.5=0.5\)，且各家相互独立，所以 \(X\) 服从二项分布 \(B\left(5,0.5\right)\). 因而
\[
P(X=2)=\mathrm C_{5}^{2}(0.5)^{2}(0.5)^{3}=\frac{10}{32}=0.3125\approx0.31.
\]

\item 由二项分布的数学期望，
\[
E(X)=5\times0.5=2.5.
\]
所以平均有 \(2.50\) 家煤矿必须整改.

\item 一家煤矿被关闭的概率为初次安检不合格且整改后复查仍不合格的概率，即 \(0.5\times(1-0.8)=0.1\). 五家煤矿均不被关闭的概率为 \((1-0.1)^{5}=0.9^{5}\)，所以至少关闭一家煤矿的概率为
\[
1-0.9^{5}=0.40951\approx0.41.
\]
\end{enumerate}
\end{solution}

\begin{problem}
如图，已知两个正四棱锥 \(P - ABCD\) 与 \(Q - ABCD\) 的高分别为 \(1\) 和 \(2\)，\(AB = 4\).

\begin{enumerate}
\item 证明：\(PQ \perp\) 平面 \(ABCD\)；
\item 求异面直线 \(AQ\) 与 \(PB\) 所成的角；
\item 求点 \(P\) 到平面 \(QAD\) 的距离.
\end{enumerate}

\FigureLayoutDeclare{img/2006/hunan_science/q18_fig1.png}{9.1cm}{14bp 19bp 22bp 12bp}
\bitmapfigure[max width=0.55\linewidth]{img/2006/hunan_science/q18_fig1.png}
\end{problem}

\begin{solution}
\begin{enumerate}
\item 设正方形 \(ABCD\) 的中心为 \(O\)，取 \(ABCD\) 所在平面为 \(z=0\)，并设 \(A(2,2,0)\)，\(B(-2,2,0)\)，\(C(-2,-2,0)\)，\(D(2,-2,0)\). 由两个正四棱锥的高及图形位置可取 \(P(0,0,1)\)，\(Q(0,0,-2)\). 因此 \(\overrightarrow{PQ}=(0,0,-3)\)，而 \(z\) 轴垂直于平面 \(ABCD\)，故 \(PQ\perp\) 平面 \(ABCD\).
\item 取 \(\overrightarrow{AQ}=(-2,-2,-2)\)，\(\overrightarrow{PB}=(-2,2,-1)\). 于是
\[
\overrightarrow{AQ}\cdot\overrightarrow{PB}=4-4+2=2,\quad \left|\overrightarrow{AQ}\right|=2\sqrt{3},\quad \left|\overrightarrow{PB}\right|=3.
\]
设异面直线 \(AQ\) 与 \(PB\) 所成的角为 \(\theta\)，则取锐角有
\[
\cos\theta=\frac{\left|\overrightarrow{AQ}\cdot\overrightarrow{PB}\right|}{\left|\overrightarrow{AQ}\right|\left|\overrightarrow{PB}\right|}=\frac{1}{3\sqrt{3}}.
\]
所以 \(\theta=\arccos\dfrac{1}{3\sqrt{3}}\).
\item \(\overrightarrow{AD}=(0,-4,0)\)，\(\overrightarrow{AQ}=(-2,-2,-2)\)，故平面 \(QAD\) 的一个法向量为
\[
\overrightarrow{AD}\times\overrightarrow{AQ}=(8,0,-8),
\]
可取法向量 \(\bs n=(1,0,-1)\). 因平面 \(QAD\) 过点 \(A(2,2,0)\)，其方程为 \(x-z-2=0\). 所以点 \(P\) 到该平面的距离为
\[
\frac{\left|0-1-2\right|}{\sqrt{1^{2}+(-1)^{2}}}=\frac{3\sqrt{2}}{2}.
\]
\end{enumerate}
\end{solution}

\begin{problem}
已知函数 \( f(x) = x - \sin x \)，数列 \(\{a_n\}\) 满足：\(0 < a_1 < 1\)，\(a_{n+1} = f(a_n)\)，\(n = 1,2,3,\cdots\). 证明：

\begin{enumerate}
\item \(0 < a_{n+1} < a_n < 1\)；
\item \(a_{n+1} < \frac{1}{6} a_n^3\).
\end{enumerate}
\end{problem}

\begin{solution}
\begin{enumerate}
\item 对 \(0<x<1\)，令 \(h(x)=x-\sin x\)，则 \(h(0)=0\)，且 \(h'(x)=1-\cos x>0\). 因而 \(x-\sin x>0\)，又因为 \(\sin x>0\)，所以 \(0<x-\sin x<x\). 由 \(0<a_1<1\) 及递推关系，用归纳法可得对任意正整数 \(n\)，都有 \(0<a_{n+1}<a_n<1\).

\item 令 \(g(x)=\sin x-x+\dfrac{x^{3}}{6}\). 有 \(g(0)=g'(0)=0\)，并且
\[
g''(x)=-\sin x+x=x-\sin x>0\quad(0<x<1).
\]
因此 \(g'(x)>0\)，进而 \(g(x)>0\)，即 \(\sin x>x-\dfrac{x^{3}}{6}\). 所以对 \(0<x<1\)，有 \(x-\sin x<\dfrac{x^{3}}{6}\). 取 \(x=a_n\)，便得到 \(a_{n+1}=a_n-\sin a_n<\dfrac{1}{6}a_n^{3}\).
\end{enumerate}
\end{solution}

\begin{problem}
对 \(1\) 个单位质量的含污物体进行清洗，清洗前其清洁度（含污物体的清洁度定义为：\(1 - \frac{\text{污物质量}}{\text{物体质量（含污物）}}\)）为 \(0.8\)，要求洗完后的清洁度是 \(0.99\). 有两种方案可供选择，方案甲：一次清洗；方案乙：两次清洗. 该物体初次清洗后受残留水等因素影响，其质量变为 \(a\)（\(1 \leqslant a \leqslant 3\)）. 设用 \(x\) 单位质量的水初次清洗后的清洁度是 \(\frac{x + 0.8}{x + 1}\)（\(x > a - 1\)），用 \(y\) 单位质量的水第二次清洗后的清洁度是 \(\frac{y + ac}{y + a}\)，其中 \(c\)（\(0.8 < c < 0.99\)）是该物体初次清洗后的清洁度.

\begin{enumerate}
\item 分别求出方案甲以及 \(c = 0.95\) 时方案乙的用水量，并比较哪一种方案用水量较少；
\item 若采用方案乙，当 \(a\) 为某定值时，如何安排初次与第二次清洗的用水量，使总用水量最少？并讨论 \(a\) 取不同数值时对最少总用水量多少的影响.
\end{enumerate}
\end{problem}

\begin{solution}
\begin{enumerate}
\item 方案甲中，设用水量为 \(x\)，则
\[
\frac{x+0.8}{x+1}=0.99,\qquad 0.01x=0.19,\qquad x=19.
\]
所以方案甲用水量为 \(19\).

方案乙在 \(c=0.95\) 时，第一次清洗用水量满足
\[
\frac{x+0.8}{x+1}=0.95,\qquad 0.05x=0.15,\qquad x=3.
\]
第二次清洗用水量 \(y\) 满足
\[
\frac{y+0.95a}{y+a}=0.99,\qquad 0.01y=0.04a,\qquad y=4a.
\]
故方案乙总用水量为 \(3+4a\). 由于 \(1\leqslant a\leqslant3\)，有 \(7\leqslant3+4a\leqslant15<19\)，所以方案乙用水量较少.

\item 设方案乙第一次清洗用水量为 \(x\)，则
\[
c=\frac{x+0.8}{x+1},\qquad y=100a(0.99-c)=\frac{a(19-x)}{x+1}.
\]
由 \(x>a-1\) 及 \(c<0.99\)，可知 \(a-1<x<19\). 总用水量为
\[
S(x)=x+\frac{a(19-x)}{x+1}.
\]
在上述区间内，
\[
S'(x)=1-\frac{20a}{(x+1)^2},\qquad S''(x)=\frac{40a}{(x+1)^3}>0.
\]
令 \(S'(x)=0\)，得唯一驻点 \(x_0=\sqrt{20a}-1\). 当 \(1\leqslant a\leqslant3\) 时，\(\sqrt{20a}>a\) 且 \(\sqrt{20a}<20\)，所以 \(a-1<x_0<19\)，驻点位于定义域内，并且由 \(S''(x)>0\) 知它给出最小值. 此时
\[
x_0=\sqrt{20a}-1,\qquad y_0=\frac{a(19-x_0)}{x_0+1}=\sqrt{20a}-a,
\]
最少总用水量为
\[
S_{\min}=x_0+y_0=2\sqrt{20a}-a-1.
\]
又
\[
\frac{\mathrm dS_{\min}}{\mathrm da}=\sqrt{\frac{20}{a}}-1>0\quad(1\leqslant a\leqslant3),
\]
所以 \(a\) 越大，最少总用水量越多. 当 \(a=1\) 时，\(S_{\min}=4\sqrt{5}-2\)；当 \(a=3\) 时，\(S_{\min}=4\sqrt{15}-4\)，故最少总用水量的取值范围为 \([4\sqrt{5}-2,\,4\sqrt{15}-4]\).
\end{enumerate}
\end{solution}

\begin{problem}
已知椭圆 \(C_{1}: \frac{x^2}{4} + \frac{y^2}{3} = 1\)，抛物线 \(C_{2}: (y - m)^2 = 2px\)（\(p > 0\)），且 \(C_{1}\)，\(C_{2}\) 的公共弦 \(AB\) 过椭圆 \(C_{1}\) 的右焦点.

\begin{enumerate}
\item 当 \(AB \perp x\) 轴时，求 \(m\)，\(p\) 的值，并判断抛物线 \(C_{2}\) 的焦点是否在直线 \(AB\) 上；
\item 是否存在 \(m\)，\(p\) 的值，使抛物线 \(C_{2}\) 的焦点恰在直线 \(AB\) 上？若存在，求出符合条件的 \(m\)，\(p\) 的值；若不存在，请说明理由.
\end{enumerate}
\end{problem}

\begin{solution}
\begin{enumerate}
\item 椭圆 \(C_1\) 的右焦点为 \(F(1,0)\). 当 \(AB\perp x\) 轴时，直线 \(AB\) 为 \(x=1\)，它与椭圆的交点为 \(\left(1,\dfrac{3}{2}\right)\) 和 \(\left(1,-\dfrac{3}{2}\right)\). 这两个点都在抛物线 \(C_2\) 上，因此
\[
\left(\frac{3}{2}-m\right)^2=2p,\qquad \left(-\frac{3}{2}-m\right)^2=2p.
\]
两式相减得 \(m=0\)，再代入得 \(2p=\dfrac{9}{4}\)，即 \(p=\dfrac{9}{8}\). 此时抛物线焦点为 \(\left(\dfrac{9}{16},0\right)\)，不在直线 \(x=1\) 上.

\item 先排除水平直线 \(AB\)：若 \(AB\) 为 \(y=0\)，它与椭圆的交点横坐标为 \(2\) 和 \(-2\)，代入抛物线分别得到 \(m^2=4p\) 和 \(m^2=-4p\)，这与 \(p>0\) 矛盾. 竖直直线的情形已在上一问中说明不满足条件. 因此设
\[
AB:\ y=k(x-1),\qquad k\neq0.
\]
将其代入椭圆方程，两个交点的横坐标 \(x_1,x_2\) 是方程
\[
(3+4k^2)x^2-8k^2x+4k^2-12=0
\]
的两个不同实根. 将直线方程代入抛物线方程，得到
\[
k^2x^2-2\bigl(k(k+m)+p\bigr)x+(k+m)^2=0.
\]
由于这两个二次方程有相同的两个根，比较对应系数可得
\[
(k+m)^2=\frac{k^2(4k^2-12)}{4k^2+3},\qquad k(k+m)+p=\frac{4k^4}{4k^2+3}.
\]
抛物线 \(C_2\) 的焦点为 \(\left(\dfrac p2,m\right)\)，它在直线 \(AB\) 上，所以 \(m=k\left(\dfrac p2-1\right)\). 令 \(r=\dfrac p2>0\)，则 \(k+m=kr\)，从而
\[
r^2=\frac{4(k^2-3)}{4k^2+3},\qquad r(k^2+2)=\frac{4k^4}{4k^2+3}.
\]
令 \(t=k^2\)，由上式知 \(t\geqslant3\)，且
\[
r=\frac{4t^2}{(4t+3)(t+2)}.
\]
将此式与 \(r^2=\dfrac{4(t-3)}{4t+3}\) 联立并平方，得
\[
4t^4=(t-3)(4t+3)(t+2)^2,
\]
即
\[
(t-6)(7t+6)(t+1)=0.
\]
结合 \(t\geqslant3\)，可知 \(t=6\). 于是 \(r=\dfrac{2}{3}\)，所以 \(p=2r=\dfrac{4}{3}\)，并且 \(m=k(r-1)=-\dfrac{k}{3}\). 由 \(k=\pm\sqrt{6}\) 得
\[
(m,p)=\left(\frac{\sqrt{6}}{3},\frac{4}{3}\right)\quad\text{或}\quad (m,p)=\left(-\frac{\sqrt{6}}{3},\frac{4}{3}\right).
\]
当 \(k=\mp\sqrt{6}\) 时，椭圆与直线有两个不同交点，且上面的系数关系保证这两个交点也在抛物线上，故上述两组值均满足条件.
\end{enumerate}
\end{solution}
